\documentclass[11pt]{article}
\usepackage[utf8]{inputenc}
\usepackage[T1]{fontenc}
\usepackage{lmodern}
\usepackage{amsmath,amssymb}
\usepackage{graphicx}
\usepackage{booktabs}
\usepackage{ulem}
\usepackage{fixltx2e}
\usepackage[margin=20mm,a4paper]{geometry}
\usepackage{hyperref}
\hypersetup{colorlinks=true,linkcolor=blue,urlcolor=blue}
\usepackage{kotex}
\title{converter}

\date{}

\begin{document}
\maketitle

\section*{NAME}\label{name}

converter - My Document Converter 의 POD 샘플

\section*{DESCRIPTION}\label{description}

POD 는 Perl 문서 형식입니다. \textbf{굵게}, \emph{기울임}, \verb|코드|,
\href{https://example.com}{링크} 를 씁니다.
This paragraph is in English.

\subsection*{목록 (Lists)}\label{목록-lists}

\begin{itemize}
\item 첫 번째 항목
\item 두 번째 항목
\item 세 번째 항목
\end{itemize}

\begin{enumerate}
\item 순서가 있는 항목
\item 두 번째
\end{enumerate}

\subsection*{코드 (Code)}\label{코드-code}

\begin{verbatim}
function greet(name) {
  console.log('Hello, ' + name)
}
\end{verbatim}

\subsection*{정의 (Definitions)}\label{정의-definitions}

\begin{description}
\item[용어] 용어에 대한 정의입니다.
\item[Pandoc] 범용 문서 변환기.
\end{description}

\end{document}
